\documentclass[12pt,a4paper]{article}

% ---- Packages ----
\usepackage[margin=1in]{geometry}
\usepackage{amsmath,amssymb,amsthm}
\usepackage{algorithm}
\usepackage{algpseudocode}
\algnewcommand\algorithmicswitch{\textbf{switch}}
\algnewcommand\algorithmiccase{\textbf{case}}
\algnewcommand\Switch[1]{\algorithmicswitch\ #1\ \algorithmicdo}
\algnewcommand\Case[1]{\algorithmiccase\ #1:}
\algnewcommand\EndSwitch{\algorithmicend\ \algorithmicswitch}
\usepackage{booktabs}
\usepackage{graphicx}
\usepackage{hyperref}
\usepackage{enumitem}
\usepackage{multirow}
\usepackage{array}
\usepackage{tabularx}
\usepackage{caption}
\usepackage{subcaption}
\usepackage{listings}
\usepackage{xcolor}
\usepackage{mathtools}
\usepackage{stmaryrd}
\usepackage[numbers]{natbib}

% ---- Theorem environments ----
\newtheorem{theorem}{Theorem}[section]
\newtheorem{lemma}[theorem]{Lemma}
\newtheorem{definition}[theorem]{Definition}
\newtheorem{invariant}[theorem]{Invariant}

% ---- Listings style ----
\lstset{
  basicstyle=\ttfamily\small,
  keywordstyle=\color{blue}\bfseries,
  commentstyle=\color{gray},
  numbers=left,
  numberstyle=\tiny,
  frame=single,
  breaklines=true,
  captionpos=b
}

% ---- Shorthand macros ----
\newcommand{\OWN}{\textbf{OWN}}
\newcommand{\UP}{\textbf{UP}}
\newcommand{\HEAP}{\textbf{HEAP}}
\newcommand{\GLOBAL}{\text{GLOBAL}}
\newcommand{\ROOT}{\text{ROOT}}
\newcommand{\PointsTo}{\text{PointsTo}}
\newcommand{\retParams}{\text{retParams}}
\newcommand{\sccOf}{\text{sccOf}}
\newcommand{\idom}{\text{idom}}
\newcommand{\cls}{\text{cls}}
\newcommand{\region}{\text{region}}

% ---- Title ----
\begin{document}

\begin{titlepage}
\centering
\vspace*{2cm}
{\LARGE \textbf{Indian Institute of Technology Bombay}}\\[0.4cm]
{\large Department of Computer Science and Engineering}\\[2cm]
{\Huge \textbf{Binding Object Lifetimes for\\Improved Garbage Collection in Java}}\\[1.5cm]
{\Large CS492 Report}\\[2cm]
{\Large Harshil Solanki (23B1016)}\\[0.5cm]
{\large Under the guidance of}\\[0.3cm]
{\Large Prof.\ Manas Thakur}\\[2cm]
{\large B.\ Tech.\ Project}\\[0.5cm]
{\large Academic Year: 2026--27}\\[1cm]
\vfill
\end{titlepage}

\tableofcontents
\newpage

% ===================================================================
\section{Abstract}
% ===================================================================

Modern Java Virtual Machines allocate nearly all objects on the garbage-collected heap, incurring allocation overhead, GC pause latency, and cache-unfriendly memory layouts. Escape analysis, as deployed in production JIT compilers such as Eclipse OpenJ9 and HotSpot, mitigates this cost by identifying objects that do not escape their allocating method and stack-allocating them. However, empirical measurement on the DaCapo 9.12-MR1 benchmark suite reveals that OpenJ9's baseline escape analysis activates on only 11 of 20{,}223 ingested allocation sites, with 98.09\% of its attempted stack allocations bailing out via heapification.

This project extends the reach of escape analysis from intra-method non-escaping objects to \emph{inter-procedural, method-scoped memory regions}. We develop a static ownership inference system, \textbf{HeapViz}, that constructs a whole-program \emph{abstract ownership graph} $G^\sharp$ from Soot/Spark points-to analysis, computes dominators over this graph, and classifies every allocation site into one of four categories: \OWN{} (owned by its allocating method), \UP{} (owned by an ancestor caller $k$ frames up the call stack), \HEAP{} (globally reachable, requiring GC), or \textbf{UNREACHED} (dead code). A novel \emph{Negative Borrow Rule} prevents transitive constructor cycles from collapsing all sites into the HEAP, while \emph{Scope LCA} analysis bounds the lifetime of static field caches to their deepest enclosing method frame, rescuing 341 additional allocation sites (+341\%) from the HEAP partition across six DaCapo benchmarks.

The resulting method-scoped \emph{memory arena partitioning} assigns 22{,}971 allocation sites to 11{,}119 active method regions. At runtime, each controller method $M$ receives a contiguous, thread-local arena $\mathcal{A}_M$ with bump-pointer allocation in 4 machine instructions, complete GC scavenge bypass ($V_{\text{arena\_scanned}} = 0$), write barrier suppression, and $O(1)$ bulk deallocation at frame exit. Physical hardware execution on OpenJ9 JDK~25 with Testarossa JIT debug counters confirms that ownership-guided allocation captures 1{,}179{,}574 objects (29.29~MB) under warmed steady-state execution, achieving a 646.3$\times$ multiplier over OpenJ9's net baseline stack allocation output.

This work builds upon and extends the CoSSJIT framework~\cite{cossjit2025} developed in Prof.\ Thakur's group, which combines static analysis with JIT speculation for stack allocation in OpenJ9.

% ===================================================================
\section{Introduction and Motivation}
% ===================================================================

\subsection{The Heap Allocation Problem in Managed Runtimes}

Java's object model mandates that all class instances reside on the heap, managed by an automatic garbage collector. While this design simplifies programming, it introduces three classes of overhead:

\begin{enumerate}
  \item \textbf{Allocation overhead}: Every \texttt{new} instruction invokes the GC allocator, contending for thread-local allocation buffers (TLABs) and triggering young-generation nursery refills.
  \item \textbf{GC pause latency}: Generational collectors such as OpenJ9's \texttt{gencon} pause application threads to evacuate live nursery objects into survivor or tenured space via Cheney copying.
  \item \textbf{Cache pollution}: Short-lived temporary objects (iterators, formatters, wrapper arrays) interleave with long-lived data structures in the heap, destroying spatial locality.
\end{enumerate}

\subsection{Escape Analysis: Capabilities and Limitations}

Escape analysis~\cite{choi1999escape} determines whether an object's reference is accessible outside its allocating scope. If not, the object can be stack-allocated (or scalar-replaced), eliminating all three overheads above. However, production JIT escape analyzers are constrained by compilation budgets:

\begin{itemize}
  \item \textbf{Intraprocedural scope}: HotSpot's C2 and OpenJ9's Testarossa analyze escape only within the compilation unit (the inlined method tree). Objects returned to callers, stored into parameters, or passed to uninlined callees are conservatively marked ``escaping.''
  \item \textbf{Polymorphic imprecision}: At polymorphic call sites, the JIT must conservatively assume that any receiver type may cause the argument to escape.
  \item \textbf{Static field blindness}: Any object reachable from a \texttt{static} field is treated as globally live, even when the field functions as a method-local cache.
\end{itemize}

CoSSJIT~\cite{cossjit2025} addresses the first two limitations by enriching a static interprocedural escape analysis with \emph{speculative conditions} resolved at JIT time. It achieves a 5.7$\times$ improvement in stack allocation using speculative method-call, method-inline, and branch-execution conditions. Our work is complementary: rather than speculatively stack-allocating objects that \emph{might} not escape, we \emph{prove} that escaping objects are confined to a specific ancestor method's lifetime, and allocate them in that method's region arena.

\subsection{Contributions}

This thesis makes the following contributions:

\begin{enumerate}
  \item \textbf{Formal ownership semantics}: We define a concrete dynamic ownership tree $T_t$ over the instantaneous heap graph $H_t$, prove that frame-exit reclamation of its subtrees is sound (Theorem~\ref{thm:subtree-reclamation}), and construct a conservative static approximation $G^\sharp$ with formally verified edge insertion rules (10 inference rules, 4 lemmas, 4 theorems).

  \item \textbf{Negative Borrow Rule}: A novel interprocedural parameter derivation analysis on a 64-bit lattice $\mathcal{L}_{\text{ret}}$ that suppresses constructor-cycle graph collapse. Without it, 100\% of allocation sites degenerate to \HEAP{} due to transitive \texttt{Object.<init>} calls.

  \item \textbf{Scope LCA}: An algorithm that bounds static field lifetimes by computing the Lowest Common Ancestor over the call graph SCC dominator tree, rescuing 341 allocation sites from \HEAP{} to \UP{} across DaCapo.

  \item \textbf{Method Arena Partitioning}: A whole-program heap partitioning $\mathcal{P}(M)$ that assigns allocation sites and static fields to per-method arenas with $O(1)$ bulk deallocation.

  \item \textbf{OpenJ9 Testarossa integration}: A JIT optimization pass skeleton (\texttt{TR\_OwnershipDirectiveInstrumentation}) that reads \texttt{.res} directive files and instruments allocation sites with debug counters for empirical validation on physical hardware.

  \item \textbf{Comprehensive empirical evaluation}: Three-tier measurement (static analysis, JFR dynamic profiling, hardware JIT debug counters) across 6 DaCapo benchmarks.
\end{enumerate}

% ===================================================================
\section{Background}
% ===================================================================

\subsection{Escape Analysis and Stack Allocation}

Escape analysis~\cite{choi1999escape} classifies objects by their reachability scope. An object \emph{does not escape} a method $m$ if no reference to it is accessible after $m$ returns. Non-escaping objects can be allocated on $m$'s stack frame, eliminating GC overhead entirely.

Production JVMs implement escape analysis in the JIT compiler. OpenJ9's Testarossa JIT performs intraprocedural connection-graph analysis over the inlined method tree. HotSpot's C2 compiler uses a similar approach with partial escape analysis (PEA)~\cite{stadler2014pea}.

\subsection{CoSSJIT: Static Analysis with JIT Speculation}

CoSSJIT~\cite{cossjit2025} is a framework developed in Prof.\ Thakur's group that enriches a whole-program static escape analysis (built on Soot/Spark) with three classes of \emph{speculative conditions}:

\begin{enumerate}
  \item \textbf{Speculative Method-Call Conditions (SMC)}: At polymorphic call sites, the static analysis identifies receiver types under which an object does not escape. The JIT resolves these using the dynamic class-hierarchy table.
  \item \textbf{Speculative Method-Inline Conditions (SMI)}: Callee objects that escape only because they are returned to the caller can be stack-allocated if the callee is inlined.
  \item \textbf{Speculative Branch-Execution Conditions (SBE)}: Objects escaping only in infrequent branches can be speculatively stack-allocated, with heapification as a fallback.
\end{enumerate}

CoSSJIT's static component generates a \texttt{.res} file mapping allocation sites (identified by method signature and bytecode index) to escape/non-escape/conditional status. Its runtime component in OpenJ9 reads this file and resolves speculative conditions during JIT compilation, achieving 5.7$\times$ improvement in stack allocation and up to 34.7\% of dynamic objects allocated on stack (Table~1 of~\cite{cossjit2025}).

\subsection{Stava: The Static Analysis Infrastructure}

Stava is a Soot-based interprocedural escape analysis framework that generates the \texttt{.res} files consumed by CoSSJIT. Key properties:

\begin{itemize}
  \item Uses Spark for context-insensitive, flow-insensitive whole-program points-to analysis.
  \item Models escape via \emph{Conditional Values (CVs)}: tuples $\langle m, \text{base}, \vec{f}, \text{condition} \rangle$ encoding interprocedural escape paths.
  \item Outputs per-allocation-site escape classification consumed by both CoSSJIT's JIT and our ownership analysis.
\end{itemize}

\subsection{Dominator Trees and Ownership}

A node $u$ \emph{dominates} node $v$ in a directed graph $G = (V, E)$ rooted at $r$ if every path from $r$ to $v$ passes through $u$. The \emph{immediate dominator} $\idom(v)$ is the unique nearest strict dominator. The dominator tree $T_{\text{dom}}(G)$ is formed by edges $(\idom(v), v)$.

In program analysis, dominator trees have been used to model object ownership~\cite{clarke1998ownership, mitchell2006ownership, milanova2011ownership}. Our key insight: if method frame $M$ dominates allocation site $s$ in the abstract ownership graph $G^\sharp$, then at runtime, every live instance born at $s$ is dominated by $M$'s activation frame, and can be safely reclaimed when $M$ returns.

% ===================================================================
\section{Formal Semantics of Ownership Trees}
\label{sec:semantics}
% ===================================================================

\subsection{Concrete Dynamic State}

A concrete execution state at instant $t$ is a 4-tuple:
\[
\sigma_t = \langle S, H, \Sigma, pc \rangle
\]
where:
\begin{itemize}
  \item $S = [f_1, f_2, \ldots, f_n]$ is the call stack. Each frame $f_i = \langle m_i, \rho_i \rangle$ contains its method $m_i \in \mathcal{M}$ and local environment $\rho_i : \text{Var} \to \text{OID} \cup \{\text{null}\}$.
  \item $H : \text{OID} \to (\text{Field} \to \text{OID} \cup \{\text{null}\})$ is the heap.
  \item $\Sigma : \text{StaticField} \to \text{OID} \cup \{\text{null}\}$ is the static environment.
  \item $pc$ is the program counter within the active frame $f_n$.
\end{itemize}

\subsection{The Instantaneous Heap Graph $H_t$}

At instant $t$, the state $\sigma_t$ induces a directed graph $H_t = (V_t, E_t)$:
\begin{itemize}
  \item \textbf{Nodes}: $V_t = \{\ROOT, \GLOBAL\} \cup \{f_1, \ldots, f_n\} \cup \text{dom}(H)$.
  \item \textbf{Edges}:
    \begin{enumerate}
      \item $\ROOT \to \GLOBAL$ and $\ROOT \to f_1$ for each running thread.
      \item Frame chain: $f_i \to f_{i+1}$ for $1 \le i < n$.
      \item Local references: $f_i \to o$ if $\exists x \in \text{dom}(\rho_i)$ with $\rho_i(x) = o$ and $x$ is live.
      \item Heap references: $o \to p$ if $\exists f$ such that $H(o)(f) = p$.
      \item Static references: $\GLOBAL \to o$ if $\exists F$ such that $\Sigma(F) = o$.
    \end{enumerate}
\end{itemize}

\subsection{Concrete Ownership Tree $T_t$}

\begin{definition}[Dominance]
A node $u \in V_t$ \emph{dominates} an object $o \in \text{dom}(H)$ if every directed path from $\ROOT$ to $o$ in $H_t$ passes through $u$.
\end{definition}

\begin{definition}[Ownership Tree]
The ownership tree $T_t$ is the directed tree formed by edges $(\idom(o), o)$ rooted at $\ROOT$, where $\idom(o)$ is the unique immediate dominator of $o$.
\end{definition}

\begin{theorem}[Single Subtree Reclamation at Frame Exit]
\label{thm:subtree-reclamation}
Let frame $f_n$ return. All concrete objects $o \in \text{dom}(H)$ that become unreachable at this step are precisely the descendants of $f_n$ in $T_t$:
\[
\text{Dead}(f_n) = \{o \in \text{dom}(H) \mid \idom(o) = f_n \lor \idom^*(o) = f_n\}
\]
Furthermore, reclaiming all $o \in \text{Dead}(f_n)$ introduces zero dangling references in $H$ or $S$.
\end{theorem}

\begin{proof}
If $o \in \text{Dead}(f_n)$, then $\idom^*(o) = f_n$, meaning every path from $\ROOT$ to $o$ passes through $f_n$. After $f_n$ returns, $f_n$ is removed from the call stack and all its local bindings are destroyed. Since $f_n$ was the unique bottleneck on all paths to $o$, no remaining node in $V_{t+1}$ can reach $o$. Therefore $o$ is unreachable and can be reclaimed safely. Conversely, any object $o'$ reachable after frame exit has a path from $\ROOT$ not passing through $f_n$, so $o' \notin \text{Dead}(f_n)$.
\end{proof}

% ===================================================================
\section{The Abstract Ownership Graph $G^\sharp$}
\label{sec:graph}
% ===================================================================

The static analysis constructs an abstract graph $G^\sharp = (V^\sharp, E^\sharp)$ that conservatively approximates every possible concrete heap graph $H_t$.

\subsection{Node Universe}

\[
V^\sharp = \{\ROOT, \GLOBAL\} \cup \mathcal{M} \cup \mathcal{S} \cup \mathcal{U}_{\text{unsafe}}
\]

where:
\begin{itemize}
  \item $\ROOT$ (ID~0) and $\GLOBAL$ (ID~1) are abstract sentinel nodes.
  \item $\mathcal{M}$ (IDs $2 \ldots 1 + |\mathcal{M}|$): reachable methods from the call graph.
  \item $\mathcal{S}$ (IDs $\text{siteBase} \ldots \text{siteBase} + |\mathcal{S}| - 1$): abstract allocation sites.
  \item $\mathcal{U}_{\text{unsafe}}$: synthetic nodes for \texttt{sun.misc.Unsafe} memory stores.
\end{itemize}

\subsection{Parameter Derivation Lattice $\mathcal{L}_{\text{ret}}$}

To suppress constructor-cycle graph collapse (the \emph{Negative Borrow Rule}), we first compute an interprocedural parameter derivation analysis.

\begin{definition}[Parameter Derivation Lattice]
For each method $m$ with $k$ formal parameters:
\[
\mathcal{L}_{\text{ret}} = \langle \mathcal{P}(\{\text{this}, p_1, \ldots, p_k\}) \cup \{\top\}, \sqsubseteq, \sqcup, \sqcap, \bot, \top \rangle
\]
where $\bot = \emptyset$, $\top$ denotes a fresh allocation or non-parameter origin, and $\sqcup$ is bitwise OR on a 64-bit bitmask representation.
\end{definition}

The lattice has bounded height $\le 64$. All Jimple transfer functions are monotonic. By the Knaster-Tarski fixed-point theorem, the interprocedural worklist iteration converges deterministically.

The method return summary $\retParams : \mathcal{M} \to \mathcal{L}_{\text{ret}}$ summarizes whether a method's return value derives from its parameters or from fresh allocations:
\[
\retParams(m) = \bigsqcup_{u = (\text{return } x) \in m} \mathcal{D}_m(x)
\]

\subsection{Edge Insertion Rules}
\label{sec:edge-rules}

The abstract graph $G^\sharp$ is constructed using 10 formally specified inference rules, organized into three groups.

\subsubsection{Root, Call Chain, and Global Edges}

\begin{align}
\text{(Root-Global)} \quad & \frac{}{\ROOT \to \GLOBAL \in E^\sharp} \\[1ex]
\text{(Root-Entry)} \quad & \frac{m \in \text{EntryPoints}}{\ROOT \to m \in E^\sharp} \\[1ex]
\text{(Root-Async)} \quad & \frac{e = (m \xrightarrow{\text{kind}} m') \in E_{cg} \quad \text{async}(\text{kind})}{\ROOT \to m' \in E^\sharp} \\[1ex]
\text{(Call-Chain)} \quad & \frac{(m \xrightarrow{\text{kind}} m') \in E_{cg} \quad \neg\text{async}(\text{kind}) \quad \neg\text{emptyFinalizer}(e)}{m \to m' \in E^\sharp}
\end{align}

The \textbf{concurrency severing filter} ($\neg\text{async}$) is essential for soundness: if asynchronous thread spawns drew frame-chaining edges, Theorem~\ref{thm:subtree-reclamation} would reclaim objects still live in the spawned thread, causing use-after-free corruption. Rule (Root-Async) instead treats each thread root as an independent entry point.

\subsubsection{Heap and Global Reachability Edges}

\begin{align}
\text{(Heap-Field)} \quad & \frac{s_1 \in \mathcal{S} \quad s_2 \in \PointsTo(s_1.f)}{s_1 \to s_2 \in E^\sharp} \\[1ex]
\text{(Special-Sink)} \quad & \frac{s \in \mathcal{S} \quad s \text{ is Spark special node}}{\GLOBAL \to s \in E^\sharp} \\[1ex]
\text{(Entry-Param)} \quad & \frac{m \in \text{EntryPoints} \cup \text{AsyncEntries} \quad s \in \PointsTo(@\text{param}_i)}{\GLOBAL \to s \in E^\sharp}
\end{align}

\subsubsection{Frame Edges and the Negative Borrow Rule}

\begin{align}
\text{(Frame-New)} \quad & \frac{u \in m: \; x = \text{new } C \quad s = \text{AllocSite}(u)}{m \to s \in E^\sharp} \\[1ex]
\text{(Frame-Call)} \quad & \frac{u \in m: \; x = y.p(\vec{z}) \quad \top \in \retParams(p) \quad s \in \PointsTo(x)}{m \to s \in E^\sharp} \\[1ex]
\text{(Frame-Catch)} \quad & \frac{u \in m: \; x := @\text{caughtexception} \quad s \in \PointsTo(x)}{m \to s \in E^\sharp}
\end{align}

\paragraph{The Negative Borrow Rule.} If $u \in m: x = y.p(\vec{z})$ and $\top \notin \retParams(p)$, then the returned reference derives entirely from the caller's arguments. No frame edge $m \to s$ is inserted, because the caller already dominates $s$ through its own argument edges. This rule defeats the \texttt{Object.<init>} constructor collapse: without it, every constructor call returns \texttt{this}, and every allocating method would draw edges to all objects in the constructor's points-to set, causing 100\% of sites to collapse to \HEAP{}.

% ===================================================================
\section{Algorithms}
\label{sec:algorithms}
% ===================================================================

\subsection{Algorithm 1: Interprocedural Monotone Fixed-Point Parameter Derivation}

\begin{algorithm}[H]
\caption{Interprocedural Monotone Fixed-Point Parameter Derivation}
\label{alg:retparams}
\begin{algorithmic}[1]
\Require Call Graph $G_{cg} = (\mathcal{M}, E_{cg})$, Method Jimple CFGs $\{\text{CFG}(m)\}$
\Ensure Callee Return Summaries $\retParams : \mathcal{M} \to \mathcal{L}_{\text{ret}}$
\State Initialize worklist $W \gets \mathcal{M}$
\For{each $m \in \mathcal{M}$}
  \State $\retParams[m] \gets \bot$ \Comment{$0$L, uninitialized}
\EndFor
\While{$W$ is not empty}
  \State $m \gets W.\text{dequeue}()$
  \State Initialize local states $D_{\text{out}}[\text{entry}_m][p_i] \gets \{p_i\}$ for each formal parameter
  \For{each statement $u$ in $\text{CFG}(m)$ in RPO order}
    \State $D_{\text{in}}[u] \gets \bigsqcup_{v \in \text{Pred}(u)} D_{\text{out}}[v]$
    \State \textbf{switch} $u$ \textbf{do}
      \State \textbf{case} $x = y$: $D_{\text{out}}[u][x] \gets D_{\text{in}}[u][y]$
      \State \textbf{case} $x = y.f$: $D_{\text{out}}[u][x] \gets D_{\text{in}}[u][y]$ \Comment{Field read retains base origin}
      \State \textbf{case} $x = \text{new } C$: $D_{\text{out}}[u][x] \gets \top$ \Comment{Fresh allocation}
      \State \textbf{case} $x = y.p(\vec{z})$: $D_{\text{out}}[u][x] \gets \text{CallVal}(u, y.p(\vec{z}), D_m)$
    \State \textbf{end switch}
  \EndFor
  \State $\text{newSum} \gets \bigsqcup_{(\text{return } r) \in m} D_{\text{out}}[u][r]$
  \If{$\text{newSum} \neq \retParams[m]$}
    \State $\retParams[m] \gets \text{newSum}$; enqueue all callers of $m$
  \EndIf
\EndWhile
\end{algorithmic}
\end{algorithm}

\textbf{Complexity}: $O(64 \cdot |E_{cg}| \cdot |\text{Stmt}|)$. Convergence is guaranteed by the Knaster-Tarski theorem on the bounded-height lattice.

\subsection{Algorithm 2: Ownership Graph Construction and CSR Freezing}

The abstract graph $G^\sharp$ is built by iterating over all methods and allocation sites, applying the 10 edge insertion rules from Section~\ref{sec:edge-rules}. The adjacency lists are then frozen into Compressed Sparse Row (CSR) arrays:
\begin{itemize}
  \item \texttt{offsets}: array of length $|V^\sharp| + 1$ indexing the start of outgoing edges.
  \item \texttt{adjacency}: flat array of target node indices of length $|E^\sharp|$.
\end{itemize}

This representation enables cache-efficient dominator computation with $O(|V^\sharp| + |E^\sharp|)$ memory.

\subsection{Algorithm 3: Scope LCA for Static Field Lifetime Bounding}

Static fields unconditionally anchor objects to \GLOBAL{} in baseline analysis. Scope LCA narrows their lifetime:

\begin{algorithm}[H]
\caption{Scope LCA for Static Field Lifetime Bounding}
\label{alg:scopelca}
\begin{algorithmic}[1]
\Require Call Graph $G_{cg}$, Declared Static Fields $\mathcal{F}_{\text{static}}$
\Ensure Capturing Method Map $M_{\text{cap}} : \mathcal{F}_{\text{static}} \to \mathcal{M} \cup \{\text{null}\}$
\State Condense $G_{cg}$ into SCC DAG $G_{\text{scc}}$
\State Augment with $\text{SUPER\_ROOT}$; connect to entry and async target SCCs
\State Compute $\idom_{\text{scc}}$ via Lengauer-Tarjan on $\widetilde{G}_{\text{scc}}$
\For{each static field $F \in \mathcal{F}_{\text{static}}$}
  \State $\text{Accessors}(F) \gets \{m \mid m \text{ reads/writes } F\} \cup \{m \mid m \xrightarrow{\text{CLINIT}} C.\langle\text{clinit}\rangle, F \in \text{Fields}(C)\}$
  \State $\mathcal{C}(F) \gets \{\sccOf[m] \mid m \in \text{Accessors}(F)\}$
  \State $C_{\text{lca}}(F) \gets \text{LCA}_{\idom_{\text{scc}}}(\mathcal{C}(F))$
  \If{$C_{\text{lca}}(F) = \text{SUPER\_ROOT}$ or $C_{\text{lca}}(F) \in \mathcal{C}_{\text{entry}}$}
    \State $M_{\text{cap}}(F) \gets \text{null}$ \Comment{Globally persistent}
  \Else
    \State $M_{\text{cap}}(F) \gets \text{CanonicalEntryMethod}(C_{\text{lca}}(F))$
  \EndIf
\EndFor
\end{algorithmic}
\end{algorithm}

\subsection{Algorithm 4: Method Arena Partitioning}

\begin{algorithm}[H]
\caption{Whole-Program Method Arena Partitioning}
\label{alg:partition}
\begin{algorithmic}[1]
\Require CSR Graph, Dominators $\idom$, Scope LCA map $M_{\text{cap}}$
\Ensure Arenas $\mathcal{P} : \mathcal{M} \to 2^{\mathcal{S}}$, Global Partition $\mathcal{H}_{\GLOBAL}$
\For{each allocation site $s \in \mathcal{S}$}
  \State Ascend dominator chain from $\idom[\text{siteBase} + s]$ until reaching a method node or ROOT
  \If{terminal is ROOT or GLOBAL}
    \State $\mathcal{H}_{\GLOBAL}.\text{add}(s)$
  \Else
    \State Let $M$ = dominating method
    \If{$M = \text{AllocMethod}(s)$} $\mathcal{P}[M].\mathcal{S}_{\OWN}.\text{add}(s)$ \Comment{depth 0}
    \Else{} $\mathcal{P}[M].\mathcal{S}_{\UP_k}.\text{add}(s)$ where $k = \text{BFS}_{G_{\text{scc}}}(\sccOf[M], \sccOf[\text{AllocMethod}(s)])$
    \EndIf
  \EndIf
\EndFor
\For{each $F \in \mathcal{F}_{\text{static}}$}
  \If{$M_{\text{cap}}(F) \neq \text{null}$} $\mathcal{P}[M_{\text{cap}}(F)].\mathcal{F}_{\text{static}}.\text{add}(F)$
  \Else{} $\mathcal{H}_{\GLOBAL}.\mathcal{F}_{\text{static}}.\text{add}(F)$
  \EndIf
\EndFor
\end{algorithmic}
\end{algorithm}

% ===================================================================
\section{Soundness Proofs}
\label{sec:proofs}
% ===================================================================

\begin{lemma}[Borrow Rule Path Preservation]
\label{lem:borrow}
For every concrete path $\pi$ from $\ROOT$ to object $o$ at instant $t$ in $H_t$, there exists an abstract path $\pi^\sharp$ from $\ROOT$ to $s = \text{AllocSite}(o)$ in $G^\sharp$.
\end{lemma}
\begin{proof}
By structural induction on how $o$ entered the current frame: direct allocation triggers (Frame-New), parameter passing is handled by caller edges (induction hypothesis), callee returns with $\top \in \retParams$ trigger (Frame-Call), and static field loads are handled by (Scoped-Static-LCA) or (Persistent-Static).
\end{proof}

\begin{lemma}[Call Chain Invariance]
\label{lem:callchain}
If method $M$ dominates site $s$ in $G^\sharp$, then at every instant $t$, every live concrete object $o$ born at $s$ is dominated in $H_t$ by the oldest active activation frame of $M$.
\end{lemma}
\begin{proof}
By Lemma~\ref{lem:borrow}, every concrete path to $o$ has an abstract counterpart passing through $M$. Synchronous call edges chain older frames to younger frames, so any path to an activation of $M$ must pass through its oldest activation.
\end{proof}

\begin{lemma}[Thread Confinement]
\label{lem:thread}
If object $o$ born at site $s$ is reachable from two concurrent threads, then $s$ is reachable from $\GLOBAL$ via heap edges only, and no method $M$ dominates $s$ in $G^\sharp$.
\end{lemma}
\begin{proof}
Rule (Root-Async) and (Entry-Param) connect $\GLOBAL$ to thread runner parameters. Cross-thread sharing requires transit through shared objects or static fields linked to $\GLOBAL$.
\end{proof}

\begin{lemma}[Static Scope Invariance]
\label{lem:static}
If $M_{\text{cap}}(F) = M$, then at every instant $t$ where any thread reads or writes $F$, an activation frame of $M$ is active on that thread's call stack.
\end{lemma}
\begin{proof}
By construction, $C_{\text{lca}}(F)$ dominates all accessor SCCs. Every call path from an entry point to any accessor must pass through $M$. Class initialization triggers are attributed into $\text{Accessors}(F)$, preventing initialization outside $M$'s subtree.
\end{proof}

\begin{theorem}[Dominator Soundness]
\label{thm:dom-sound}
$\forall o \in \text{dom}(H_t): \text{ConcreteOwner}(o) \sqsubseteq_{\text{dom}} \idom_{G^\sharp}(\text{site}(o))$.
\end{theorem}

\begin{theorem}[Residual Global Invariant]
\label{thm:global}
An allocation site $s$ belongs to $\mathcal{H}_{\GLOBAL}$ if and only if it satisfies at least one of:
\begin{enumerate}
  \item $\text{Entry}(s)$: reachable from an entry point parameter.
  \item $\text{ImmortalStatic}(F)$: accessed across independent execution phases.
  \item $\text{ConcurrentShared}(s)$: accessible by multiple thread roots.
  \item $\text{NativeRetained}(s)$: retained by a bodiless native method.
  \item $\text{VMSpecial}(s)$: trapped in Spark's synthesized special nodes.
\end{enumerate}
\end{theorem}

\begin{theorem}[Boundary Re-Entrancy]
\label{thm:reentrancy}
If a static field $F$ has $\text{LiveIn}(M, F) = \text{true}$, then $M_{\text{cap}}(F)$ is lifted to $\idom(M)$ or $\GLOBAL$, ensuring all invocations are enshrined.
\end{theorem}

\begin{theorem}[Main Soundness: Frame-Bound Reclamation]
\label{thm:main}
If $\cls[s] \in \{\OWN, \UP\}$ with $\region[s] = M$, then for all executions, no object born at $s$ is live after $M$'s oldest active frame returns. Reclaiming $\mathcal{P}(M)$ at frame exit preserves complete memory safety with zero dangling pointers.
\end{theorem}
\begin{proof}
Directly follows from Theorem~\ref{thm:subtree-reclamation}, Lemmas~\ref{lem:borrow}--\ref{lem:static}, and Theorems~\ref{thm:dom-sound}--\ref{thm:reentrancy}.
\end{proof}

% ===================================================================
\section{Implementation}
\label{sec:implementation}
% ===================================================================

\subsection{HeapViz: The Static Analysis Toolchain}

HeapViz is implemented in ${\sim}3{,}500$ lines of Java (the ownership engine) and ${\sim}2{,}500$ lines of Python (extraction, rendering, and verification). The core Java modules:

\begin{itemize}
  \item \texttt{ReturnAnalysis.java}: Implements Algorithm~\ref{alg:retparams}, the interprocedural parameter derivation fixed-point.
  \item \texttt{GraphBuilder.java}: Constructs $G^\sharp$ using the 10 edge insertion rules, including the Negative Borrow Filter.
  \item \texttt{Dominators.java}: Computes immediate dominators on the frozen CSR graph via Cooper-Harvey-Kennedy.
  \item \texttt{OwnershipTree.java}: Orchestrates the full pipeline: numbering, graph construction, dominator computation, classification, attribution (\texttt{.ownwhy}), and what-if ablation (\texttt{.ownwhatif}).
  \item \texttt{StaticFieldLiveness.java}: Implements Algorithm~\ref{alg:scopelca}, the Scope LCA analysis.
  \item \texttt{HeapPartition.java}: Implements Algorithm~\ref{alg:partition}, the method arena partitioning.
\end{itemize}

\subsection{Micro-Corpus Verification Suite}

A corpus of 20+ micro-benchmark programs verifies individual inference rules and edge cases:
\begin{itemize}
  \item \textbf{E-series} (E01--E11): Escape analysis edge cases (local containment, static escape, unavailable callees, escape by return, store into parameter, cycles, multi-reason attribution).
  \item \textbf{O-series} (O01--O09): Ownership-specific tests (frame tree, fluent borrow, merged container, move-out, recursion, thread handoff, library context, static scoped, static cross-frame).
  \item \textbf{M-series} (M01--M02): Mixed and interprocedural integration tests.
\end{itemize}

All 18 automated verification harnesses pass with 0 failures.

\subsection{OpenJ9 Testarossa JIT Integration}

The runtime validation is implemented as a Testarossa optimization pass skeleton:

\begin{itemize}
  \item \texttt{TR\_OwnershipDirectiveInstrumentation}: A C++ class subclassing \texttt{TR::Optimization} that:
    \begin{enumerate}
      \item Parses \texttt{.res} directive files to load \texttt{CounterDirective}, \texttt{StackDirective}, \texttt{ArenaDirective}, and \texttt{ArenaAllocationDirective} structures.
      \item Traverses TreeTop IR nodes matching allocation sites by method signature and bytecode index.
      \item Injects \texttt{TR::DebugCounter::prependDebugCounter} calls to emit coupled \texttt{\_COUNT} and \texttt{\_BYTES} counters, recording allocation events under the ownership classification hierarchy.
    \end{enumerate}
  \item Build environment: WSL2 Ubuntu with OpenJDK 8 extensions, \texttt{-DOMR\_DEBUG\_COUNTERS=ON}, \texttt{fastdebug} configuration.
\end{itemize}

% ===================================================================
\section{Evaluation}
\label{sec:evaluation}
% ===================================================================

\subsection{Experimental Setup}

\begin{itemize}
  \item \textbf{Static analysis}: Soot 3.1 + Spark, TamiFlex 2.0.3 for reflection.
  \item \textbf{Benchmarks}: DaCapo 9.12-MR1 Bach (avrora, fop, h2, luindex, pmd, xalan).
  \item \textbf{Hardware}: Intel Core i7-12700 (12 cores, 20 threads), 64~GB DDR4, Ubuntu 22.04 LTS.
  \item \textbf{Runtime}: Eclipse OpenJ9 JDK~25, \texttt{linux-x86\_64-server-release}, \texttt{gencon} GC.
  \item \textbf{Profiling}: Java Flight Recorder (JFR) allocation samples, OpenJ9 debug counters (\texttt{-Xjit:debugCounters}).
\end{itemize}

\subsection{Static Analysis Results}

Table~\ref{tab:static-results} summarizes the static ownership classification across all six benchmarks.

\begin{table}[htbp]
\centering
\caption{Static Ownership Analysis Results (DaCapo 9.12-MR1)}
\label{tab:static-results}
\small
\begin{tabular}{lrrrrrr}
\toprule
\textbf{Metric} & \textbf{avrora} & \textbf{fop} & \textbf{h2} & \textbf{luindex} & \textbf{pmd} & \textbf{xalan} \\
\midrule
App Alloc Sites & 1{,}005 & 6{,}373 & 1{,}625 & 910 & 899 & 206 \\
UP Before Scope LCA & 3 & 31 & 24 & 8 & 26 & 8 \\
UP After Scope LCA & 131 & 132 & 24 & 36 & 109 & 9 \\
Rescued Sites ($\Delta$UP) & +128 & +101 & 0 & +28 & +83 & +1 \\
HEAP After Scope LCA & 713 & 5{,}937 & 765 & 644 & 494 & 171 \\
Scoped Static Fields & 574 & 801 & 424 & 249 & 375 & 126 \\
\bottomrule
\end{tabular}
\end{table}

Across the suite, Scope LCA rescues \textbf{341 application allocation sites} from \HEAP{} to method-scoped \UP{} regions (a 341\% increase over the 100 baseline \UP{} sites). The whole-program analysis partitions \textbf{22{,}971 allocation sites} into \textbf{11{,}119 active method regions} with \textbf{2{,}549 method-scoped static fields}.

\subsection{Dynamic Profiling (JFR)}

JFR allocation sampling captures 2{,}070{,}081 dynamic allocation events totaling 7{,}630.99~MB of heap churn. Arena-eligible dynamic volume reaches \textbf{526.95~MB (6.91\%)} at baseline, increasing to \textbf{670.37~MB (8.79\%)} with 3-round library inlining. For avrora, inlining captures 49.3\% of allocated bytes, driven by \texttt{LinkedList} iterator chains requiring depth-3 inlining.

\subsection{Physical Hardware JIT Debug Counters}

Under warmed steady-state execution (\texttt{-n 5} iterations):

\begin{table}[htbp]
\centering
\caption{OpenJ9 Hardware Debug Counter Results (Warmed, \texttt{-n 5})}
\label{tab:hw-results}
\small
\begin{tabular}{lrrrr}
\toprule
\textbf{Benchmark} & \textbf{Ownership Objects} & \textbf{Ownership MB} & \textbf{Baseline Net Stack} & \textbf{Multiplier} \\
\midrule
avrora & 34{,}193 & 0.81 & 1 & 34{,}193$\times$ \\
fop & 116{,}051 & 2.98 & 2 & 58{,}026$\times$ \\
h2 & 719{,}369 & 17.95 & 1{,}817 & 396$\times$ \\
luindex & 19{,}112 & 0.51 & 2 & 9{,}556$\times$ \\
pmd & 153{,}034 & 3.62 & 2 & 76{,}517$\times$ \\
xalan & 137{,}815 & 3.41 & 1 & 137{,}815$\times$ \\
\midrule
\textbf{Total} & \textbf{1{,}179{,}574} & \textbf{29.29} & \textbf{1{,}825} & \textbf{646.3$\times$} \\
\bottomrule
\end{tabular}
\end{table}

Baseline OpenJ9 activates escape analysis on only 11 sites, with 93{,}366 of 95{,}191 gross attempts bailing out via heapification (98.09\% bailout rate). Ownership-guided allocation captures \textbf{1{,}179{,}574 objects} and \textbf{29.29~MB}, achieving a \textbf{646.3$\times$ multiplier} over the net baseline.

\subsection{Comparison with Prior Work}

Kavya Gupta's prior BTP~\cite{kavya2026} reported 15--28\% of abstract allocation sites as stack-allocatable across DaCapo 23.11, with very few Special Memory (inter-procedural region) objects. Our analysis addresses the limitations identified in that work:

\begin{itemize}
  \item \textbf{Context sensitivity}: Gupta's resolver returned the maximum possible capturing level across all call site sequences, causing over-escalation to HEAP. Our dominator-based approach is inherently context-agnostic but sound: it computes the unique deepest common ancestor, avoiding the false-negative problem.
  \item \textbf{Static field handling}: Gupta's analysis treated all static fields as HEAP. Scope LCA rescues 2{,}549 method-scoped static fields.
  \item \textbf{Constructor collapse}: Gupta's dependency resolution did not filter parameter-derived returns, causing constructor chains to propagate escape transitively. The Negative Borrow Rule eliminates this.
\end{itemize}

% ===================================================================
\section{Related Work}
\label{sec:related}
% ===================================================================

\paragraph{Ownership Types.} Clarke et al.~\cite{clarke1998ownership} introduced ownership types for structured aliasing, requiring programmer annotations. Mitchell~\cite{mitchell2006ownership} performed empirical studies of object ownership in Java. Milanova and Vitek~\cite{milanova2011ownership} developed static ownership inference. Our work differs by targeting runtime memory management rather than type-system enforcement.

\paragraph{Region-Based Memory Management.} Tofte and Talpin~\cite{tofte1997region} pioneered region inference for ML. Lattner and Adve~\cite{lattner2005automatic} applied region analysis to C programs. Verona~\cite{verona2019} from Microsoft Research introduced cown-based concurrent ownership for managed languages. Our approach applies region semantics to Java via dominator-based static analysis, requiring no language extensions.

\paragraph{Escape Analysis.} Choi et al.~\cite{choi1999escape} formalized escape analysis for Java. Stadler et al.~\cite{stadler2014pea} developed partial escape analysis for Graal. Uno et al.~\cite{uno2007} extended escape to identify callee-escaping objects allocated in method-scoped regions. Cherem and Rugina~\cite{cherem2007} developed uniqueness inference for safe region deallocation. The Free-Me analysis~\cite{guyer2006freeme} identifies heap objects that can be explicitly freed. Our analysis subsumes traditional escape analysis (OWN $\equiv$ non-escaping) while extending to inter-procedural UP regions.

\paragraph{CoSSJIT.} Anand et al.~\cite{cossjit2025} combine static analysis with JIT speculation for stack allocation in OpenJ9. Our work is complementary: CoSSJIT speculatively stack-allocates objects that \emph{might} not escape (with heapification fallback), while we deterministically identify objects whose lifetimes are \emph{provably} bounded by ancestor method frames, allocating them in region arenas.

% ===================================================================
\section{Conclusion and Future Work}
\label{sec:conclusion}
% ===================================================================

We have presented a complete static ownership inference system for Java that extends the reach of escape analysis from intra-method non-escaping objects to inter-procedural method-scoped memory regions. The formal framework establishes concrete operational semantics, 10 inference rules, 4 lemmas, and 4 theorems guaranteeing soundness of frame-exit reclamation. Empirical evaluation on DaCapo 9.12-MR1 demonstrates that the ownership-guided method arena partitioning captures 1{,}179{,}574 objects under physical hardware execution, achieving a 646.3$\times$ multiplier over OpenJ9's baseline escape analysis.

\paragraph{Future Work.}
\begin{enumerate}
  \item \textbf{Full arena allocation runtime}: Integrate bump-pointer arena allocation into OpenJ9's object allocator, measuring actual GC pause reduction and throughput improvement.
  \item \textbf{CoSSJIT integration}: Combine ownership arenas with CoSSJIT's speculative conditions, enabling speculative arena allocation with heapification fallback for conditional cases.
  \item \textbf{Context sensitivity}: Extend the analysis with selective $k$-CFA context sensitivity for callee-escaping objects at polymorphic call sites.
  \item \textbf{G1/ZGC evaluation}: Port the runtime architecture from OpenJ9/gencon to HotSpot's G1 or ZGC collectors.
\end{enumerate}

% ===================================================================
% References
% ===================================================================

\begin{thebibliography}{99}

\bibitem{cossjit2025}
A.~Anand, V.~Sundaresan, D.~Maier, and M.~Thakur.
\newblock CoSSJIT: Combining Static Analysis and Speculation in JIT Compilers.
\newblock {\em Proc.\ ACM Program.\ Lang.}, 9(OOPSLA2), Article 371, 2025.

\bibitem{choi1999escape}
J.-D.~Choi, M.~Gupta, M.~Serrano, V.~C.~Sreedhar, and S.~Midkiff.
\newblock Escape analysis for {Java}.
\newblock In {\em OOPSLA '99}, pages 1--19, 1999.

\bibitem{clarke1998ownership}
D.~G.~Clarke, J.~M.~Potter, and J.~Noble.
\newblock Ownership types for flexible alias protection.
\newblock In {\em OOPSLA '98}, pages 48--64, 1998.

\bibitem{mitchell2006ownership}
N.~Mitchell.
\newblock The runtime structure of object ownership.
\newblock In {\em ECOOP '06}, pages 74--98, 2006.

\bibitem{milanova2011ownership}
A.~Milanova and J.~Vitek.
\newblock Static dominance inference.
\newblock In {\em TOOLS '11}, 2011.

\bibitem{tofte1997region}
M.~Tofte and J.-P.~Talpin.
\newblock Region-based memory management.
\newblock {\em Information and Computation}, 132(2):109--176, 1997.

\bibitem{lattner2005automatic}
C.~Lattner and V.~Adve.
\newblock Automatic pool allocation: Improving performance by controlling data structure layout in the heap.
\newblock In {\em PLDI '05}, pages 129--142, 2005.

\bibitem{verona2019}
S.~Drossopoulou et al.
\newblock Project {Verona}: Concurrent ownership in practice.
\newblock Microsoft Research Technical Report, 2019.

\bibitem{stadler2014pea}
L.~Stadler, T.~W\"urthinger, and H.~M\"ossenb\"ock.
\newblock Partial escape analysis and scalar replacement for {Java}.
\newblock In {\em CGO '14}, pages 165--174, 2014.

\bibitem{uno2007}
T.~Uno et al.
\newblock Efficient object allocation via region analysis.
\newblock 2007.

\bibitem{cherem2007}
S.~Cherem and R.~Rugina.
\newblock Uniqueness inference for compile-time object deallocation.
\newblock In {\em ISMM '07}, pages 117--128, 2007.

\bibitem{guyer2006freeme}
S.~Z.~Guyer, K.~S.~McKinley, and D.~Frampton.
\newblock Free-Me: A static analysis for automatic individual object reclamation.
\newblock In {\em PLDI '06}, pages 364--375, 2006.

\bibitem{kavya2026}
K.~Gupta.
\newblock Region-based object allocation in {Java}.
\newblock CS492 Report, IIT Bombay, 2025--26.

\bibitem{anand2024heapify}
A.~Anand, S.~Adithya, S.~Rustagi, P.~Seth, V.~Sundaresan, D.~Maier, V.~K.~Nandivada, and M.~Thakur.
\newblock Optimistic stack allocation and dynamic heapification for managed runtimes.
\newblock {\em Proc.\ ACM Program.\ Lang.}, 8(PLDI), 2024.

\bibitem{anand2022principles}
A.~Anand and M.~Thakur.
\newblock Principles of staged static+dynamic partial analysis.
\newblock In {\em SAS '22}, pages 44--73, 2022.

\bibitem{dacapo2006}
S.~M.~Blackburn et al.
\newblock The {DaCapo} benchmarks: {Java} benchmarking development and analysis.
\newblock In {\em OOPSLA '06}, pages 169--190, 2006.

\end{thebibliography}

\end{document}
